\documentclass[11pt,a4paper]{article}
\usepackage{amsmath,amssymb}
\usepackage[margin=2.5cm]{geometry}
\usepackage{enumitem}
\usepackage{hyperref}
\usepackage{fancyhdr}
\usepackage{microtype}

\pagestyle{fancy}
\fancyhf{}
\rhead{ETH Zurich Executive Education}
\lhead{Section 5 --- Annotated Bibliography}
\cfoot{\thepage}

\title{How Close Are We to AGI? \\[0.2em]
       \large Annotated Bibliography for the ETH Zurich Executive Lecture}
\author{}
\date{}

\begin{document}
\maketitle
\thispagestyle{fancy}

\noindent
Six anchor references shape this lecture. A further cluster of supporting
references is provided for participants who wish to read beyond the session.
Each entry carries a one-sentence plain-English takeaway suitable for a
manager's reading-list briefing.

\section*{Six Anchor References}

\subsection*{A1. Morris et al.\ (DeepMind), \emph{Levels of AGI} (2024)}

\textbf{Full citation.}
Morris, M.\,R., Sohl-Dickstein, J., Fiedel, N., Warkentin, T., Dafoe, A.,
Faust, A., Farabet, C.\ and Legg, S.\
\emph{Levels of AGI for Operationalizing Progress on the Path to AGI}.
arXiv:2311.02462 (2024).
Published as a Position paper at ICML 2024.

\textbf{What it says.}
Proposes a 5$\times$2 matrix: five levels of performance (Emerging,
Competent, Expert, Virtuoso, Superhuman) crossed with two levels of
generality (Narrow, General). Today's frontier models are placed at
\emph{Emerging AGI} (General).

\textbf{Why it matters for managers.}
Before arguing about ``when AGI arrives,'' you need a shared definition.
This is the best-engineered working definition available, authored by
the research group (DeepMind) most likely to build one.

\subsection*{A2. Bubeck et al.\ (Microsoft), \emph{Sparks of AGI} (2023)}

\textbf{Full citation.}
Bubeck, S., Chandrasekaran, V., Eldan, R., Gehrke, J., Horvitz, E.,
Kamar, E., Lee, P., Lee, Y.\,T., Li, Y., Lundberg, S., Nori, H., Palangi,
H., Ribeiro, M.\,T.\ and Zhang, Y.
\emph{Sparks of Artificial General Intelligence: Early experiments
with GPT-4}. arXiv:2303.12712 (2023).

\textbf{What it says.}
Microsoft researchers with early GPT-4 access document qualitative
demonstrations of problem-solving across mathematics, code, vision,
medicine, law and theory-of-mind tasks. Argue the behaviour is
``reasonably viewed as an early (yet still incomplete) version of
an artificial general intelligence system.''

\textbf{Why it matters for managers.}
This is the paper that moved the AGI debate from internet forums into
the board room. Reading the demonstrations and the authors' caveats
is essential; taking the paper's headline claim at face value is not
recommended.

\subsection*{A3. Trinh et al.\ (DeepMind), \emph{AlphaGeometry} (2024)}

\textbf{Full citation.}
Trinh, T.\,H., Wu, Y., Le, Q.\,V., He, H.\ and Luong, T.
\emph{Solving olympiad geometry without human demonstrations}.
\emph{Nature} \textbf{625}, 476--482 (2024).

\textbf{What it says.}
A neuro-symbolic system (language model + symbolic deduction engine)
solves 25 of 30 IMO-level geometry problems --- comparable to a gold
medallist (25.9 average). Training uses 100\,million synthetic theorems
generated algorithmically, with no human proofs.

\textbf{Why it matters for managers.}
Proof-of-principle that (a) AI can now do \emph{creative} mathematical
reasoning, and (b) in domains with a formal verifier, human data is no
longer a bottleneck. Follow-on work (AlphaProof, AlphaGeometry 2, 2024)
achieved silver medal at the IMO itself. Expect any domain with a formal
verifier to fall fast.

\subsection*{A4. METR, \emph{Measuring the Ability of Autonomous Agents to
Complete Tasks} (2024--2026)}

\textbf{Full citation.}
Model Evaluation and Threat Research (METR). Technical reports and
blog posts, 2024--2026. \url{https://metr.org/blog}

\textbf{What it says.}
Defines a task-horizon metric: the time-length of a task on which an
autonomous agent achieves 50\% success. Tracks how this horizon has
scaled with model generation. Reports a doubling approximately every
seven months over the recent window.

\textbf{Why it matters for managers.}
If the trend continues, multi-day autonomous task completion arrives
in the late 2020s. Whether or not the extrapolation holds, the METR
trajectory is the most business-relevant capability signal currently
published; tracking it quarterly is a reasonable policy.

\subsection*{A5. Grace et al., \emph{Thousands of AI Authors on the Future of
AI} (2024)}

\textbf{Full citation.}
Grace, K., Stewart, H., Sandkuehler, J.\,F., Thomas, S., Weinstein-Raun,
B.\ and Brauner, J.
\emph{Thousands of AI Authors on the Future of AI}.
arXiv:2401.02843 (2024).

\textbf{What it says.}
Largest expert survey to date (2\,778 respondents, all published authors
at top AI venues). Aggregate 50\% forecast for High-Level Machine
Intelligence: 2047 --- 13 years earlier than the same survey gave in
2022. 38\% of respondents assign at least 10\% probability to
``extremely bad outcomes'' from advanced AI.

\textbf{Why it matters for managers.}
Expert forecasts are not ground truth, but this is the largest and
most rigorously gathered sample available. The key insight: the
experts disagree by decades --- that disagreement is itself a signal
about the state of knowledge.

\subsection*{A6. Bengio, Hinton, Russell et al., \emph{Managing Extreme AI
Risks Amid Rapid Progress} (\emph{Science}, 2024)}

\textbf{Full citation.}
Bengio, Y., Hinton, G., Yao, A., Song, D., Abbeel, P., Darrell, T.,
Harari, Y.\,N., Zhang, Y.-Q., Xue, L., Shalev-Shwartz, S., Hadfield, G.,
Clune, J., Maharaj, T., Hutter, F., Baydin, A.\,G., McIlraith, S., Gao,
Q., Acharya, A., Krueger, D., Dragan, A., Torr, P., Russell, S.,
Kahneman, D., Brauner, J.\ and Mindermann, S.
\emph{Managing extreme AI risks amid rapid progress}.
\emph{Science} \textbf{384}(6698), 842--845 (2024).
DOI: 10.1126/science.adn0117.

\textbf{What it says.}
A consensus statement by leading AI researchers, including two Turing
Award winners (Hinton, Yao) and a Nobel-Prize recipient (Kahneman),
arguing that current AI trajectories create plausible existential-scale
risks. Concrete governance asks include mandatory pre-deployment
evaluation, national AI safety institutes, compute-based licensing for
frontier training runs, and developer liability.

\textbf{Why it matters for managers.}
When Nobel and Turing laureates publish a joint statement in
\emph{Science}, the appropriate managerial response is not to decide
``who's right'' but to ensure the policy infrastructure, internal
governance and evaluation capability exist to respond to whichever
scenario plays out.

\section*{Supporting References}

\subsection*{S1. Kaplan et al., \emph{Scaling Laws for Neural Language Models}
(2020)}
OpenAI's foundational scaling paper: bigger models + more data + more
compute yield predictable, smooth improvements in loss. The paper that
justified the gigascale race.

\subsection*{S2. Hoffmann et al., \emph{Training Compute-Optimal Large
Language Models} (``Chinchilla,'' NeurIPS 2022)}
DeepMind's correction to Kaplan et al.: current models are undertrained
relative to their size; the optimal recipe uses many more tokens per
parameter. Rewrote the scaling playbook.

\subsection*{S3. Wei et al., \emph{Emergent Abilities of Large Language Models}
(TMLR 2022)}
Showed that some capabilities ``jump'' sharply at a critical scale,
rather than scaling smoothly. The paper behind the ``emergence'' headline.

\subsection*{S4. Schaeffer, Miranda, Koyejo, \emph{Are Emergent Abilities a
Mirage?} (NeurIPS 2023, Outstanding Paper)}
Counter-result: many emergent jumps disappear when the metric is
continuous rather than thresholded. The emergence is largely an artefact
of measurement, not capability.

\subsection*{S5. Chollet, \emph{On the Measure of Intelligence} (2019)}
The strongest skeptical framework. Argues that scale-driven improvements
are on \emph{skill}, not \emph{intelligence}, and proposes ARC-AGI as a
benchmark designed to resist this confusion.

\subsection*{S6. Yao et al., \emph{ReAct} (ICLR 2023)}
The template for modern LLM agents: interleave reasoning steps with
action/tool-use steps. The architectural ancestor of most current agent
frameworks.

\subsection*{S7. Shinn et al., \emph{Reflexion} (NeurIPS 2023)}
Agents that learn from their own failures by generating verbal
self-critique, without updating weights. The pattern that made
``iterative'' agents work.

\subsection*{S8. Park et al., \emph{Generative Agents} (UIST 2023, Best Paper)}
A populated simulation of LLM-based agents forming memories, plans,
relationships and collective behaviour. The paper that brought
multi-agent LLM systems into mainstream view.

\subsection*{S9. Jimenez et al., \emph{SWE-bench} (ICLR 2024)}
The most credible ``real work'' benchmark for LLM agents: can the
system fix real GitHub issues in real open-source projects? Progress
on SWE-bench (and SWE-bench Verified) is the standard agent-progress
meter.

\subsection*{S10. Mialon et al., \emph{GAIA} (ICLR 2024)}
A benchmark designed for general AI assistants, intentionally hard for
current frontier models (even GPT-4 scored 15\% at launch on the
hardest tier), easy for humans.

\subsection*{S11. FAIR Diplomacy Team, \emph{CICERO} (\emph{Science} 2022)}
Human-level play in Diplomacy by fusing language models with
game-theoretic strategic reasoning. The first serious demonstration of
language-plus-strategy AI.

\subsection*{S12. Bommasani et al., \emph{On the Opportunities and Risks of
Foundation Models} (Stanford CRFM, 2021)}
The report that introduced and defined the term ``foundation model.''
A reference document for how the field now talks about itself.

\subsection*{S13. Hendrycks, Mazeika, Woodside, \emph{An Overview of
Catastrophic AI Risks} (2023)}
Structured taxonomy of AI risk classes: malicious use, AI race, systemic
risks, and rogue AI. Complements the Bengio et al.\ Science paper.

\subsection*{S14. Stanford AI Index Report 2024 (HAI)}
The definitive annual reference. The chart-book for anyone briefing a
board on the state of AI.

\subsection*{S15. Epoch AI, \emph{Compute Trend Tracker} (2024--2026)}
Data that grounds scaling arguments. The public record of training
compute, model size, and training tokens across frontier AI systems.

\subsection*{S16. UK AI Safety Institute, \emph{Approach to Evaluations} (2024)}
The first government-grade evaluation methodology for frontier AI.
Worth reading to see what an accountable pre-deployment-evaluation
regime actually looks like.

\subsection*{S17. Silver, Singh, Precup, Sutton, \emph{Reward is Enough}
(\emph{AI Journal} 2021)}
The provocative thesis that reward maximisation in a sufficiently rich
environment is itself enough to produce general intelligence. Reading
list candidate for the ambitious.

\section*{How To Use This Bibliography}

\begin{itemize}
    \item \textbf{If you read nothing else:} Morris et al.\ (A1) and Bengio
          et al.\ (A6). Five minutes and a tea.
    \item \textbf{If you have an afternoon:} add AlphaGeometry (A3) for the
          ``wait, what?'' moment, Grace et al.\ (A5) for forecasting, and
          the Stanford AI Index (S14) for breadth.
    \item \textbf{If you want to argue with confidence:} add Schaeffer
          et al.\ (S4) and Chollet (S5) --- the two strongest sceptical
          positions currently in the literature.
    \item \textbf{If you want to understand what AI agents actually do:}
          Yao et al.\ (S6), Shinn et al.\ (S7), Jimenez et al.\ (S9) and
          the METR reports (A4).
\end{itemize}

\end{document}
